\section{HPC Case Study: Photonic Acceleration for Graph Partitioning}
\label{sec:hpc}

\subsection{Motivation}

Graph partitioning is a critical preprocessing step in high-performance computing (HPC) for domain decomposition, load balancing, and communication minimization in parallel simulations. The Minimum Bisection problem arises directly in:
\begin{itemize}
\item \textbf{Finite element mesh partitioning:} Balancing elements across MPI ranks while minimizing interface communication.
\item \textbf{Sparse matrix partitioning:} Balancing rows/columns for parallel sparse linear algebra (e.g., PETSc, Trilinos).
\item \textbf{Graph neural network training:} Balanced partitioning for distributed GNN training (e.g., DistDGL, PyTorch Distributed).
\end{itemize}

Current HPC practice uses multilevel algorithms (METIS, KaHIP, Scotch) which are fast but heuristic. For small-to-medium graphs ($n \le 100$), exact or near-exact partitioning can significantly reduce communication volume.

\subsection{Photonic Acceleration Potential}

A photonic QPU could serve as a \textbf{specialized co-processor} for graph partitioning:
\begin{itemize}
\item \textbf{Latency:} Optical interferometry operates at speed of light; a single shot takes $\sim$ns.
\item \textbf{Parallelism:} All $n$ modes evolve simultaneously; the interferometer computes the full output distribution in one shot.
\item \textbf{Energy efficiency:} Passive linear optics consumes no energy for the computation itself (only for photon generation and detection).
\end{itemize}

\subsection{Integration Workflow}

\begin{enumerate}
\item \textbf{Preprocessing (CPU):} Graph coarsening to $n \le 12$ (or $n \le 20$ with future devices).
\item \textbf{Photonic kernel (QPU):} Variational optimization + sampling to obtain high-quality balanced partitions.
\item \textbf{Postprocessing (CPU):} Uncoarsening and local refinement (Kernighan-Lin, FM).
\end{enumerate}

\subsection{Projected Performance}

For a graph with $n=100$ vertices, coarsened to $n=12$:
\begin{itemize}
\item Classical METIS: $\sim$1 ms, cut quality within 5\% of optimum.
\item Photonic kernel (projected): 10,000 shots at 1 MHz rep rate = 10 ms, expected $P_{\mathrm{success}} \approx 0.03$ (from Table~\ref{tab:variational_main}), so $\sim$300 shots for one optimal solution = 0.3 ms.
\item Total: $\sim$1.3 ms with potential for better cut quality.
\end{itemize}

The photonic kernel is not a replacement for METIS but a \textbf{high-quality initializer} that can find better partitions for difficult graphs where multilevel heuristics get stuck in local minima.

\subsection{Challenges for HPC Deployment}

\begin{itemize}
\item \textbf{Graph size:} Current QPUs support $n \le 12$; coarsening adds overhead.
\item \textbf{Repetition rate:} Current single-photon sources $\sim$MHz; need $\sim$GHz for competitive latency.
\item \textbf{Deterministic sources:} Heralded sources have probabilistic emission; deterministic quantum dot sources needed.
\item \textbf{Integration:} Optical I/O bandwidth and CPU-QPU communication latency.
\end{itemize}

\subsection{Conclusion}

Photonic quantum optimization is not a universal HPC accelerator but a \textbf{specialized tool for high-quality graph partitioning} where cut quality is paramount and graph size permits coarsening to QPU scale. The variational collision-aware approach developed here provides the algorithmic foundation; hardware advances in deterministic sources and high-rep-rate detection will determine practical viability.
